% Eigenständiger experimenteller Lesetext. Keine Einbindung in die Satzregistries.
% Die Randnotation gehört bewusst nur zu dieser Ausgabe.
\documentclass[a4paper,11pt,oneside]{report}
\usepackage{fontspec}
\setmainfont{Latin Modern Roman}
\usepackage[ngerman]{babel}
\usepackage{csquotes}
\usepackage{amsmath,amssymb,mathtools}
\usepackage[a4paper,left=24mm,right=47mm,top=24mm,bottom=25mm,
  marginparwidth=31mm,marginparsep=5mm]{geometry}
\usepackage{microtype}
\usepackage{booktabs,tabularx}
\usepackage{needspace}
\usepackage{tcolorbox}
\tcbuselibrary{breakable}
\usepackage[hidelinks,unicode]{hyperref}
\hypersetup{pdftitle={Was zweiseitige Produkte verraten - Experimentelle Lesefassung},
  pdfauthor={Martin Kunze}}
\normalmarginpar
\setlength{\parindent}{0pt}
\setlength{\parskip}{.55\baselineskip}
\setlength{\marginparpush}{12pt}
\newcommand{\Rand}[2]{\marginpar{\raggedright\footnotesize
  \textbf{\boldmath #1}\par\smallskip #2}}
\newcommand{\Abschnitt}[2]{\par\Needspace{6\baselineskip}\section*{#1}
  \phantomsection\label{#2}\addcontentsline{toc}{section}{#1}}
\newtcolorbox{Lesebeispiel}[1]{colback=white,colframe=black!45,
  colbacktitle=black!5,coltitle=black,fonttitle=\bfseries,title={#1},
  boxrule=.4pt,arc=0pt,left=7pt,right=7pt,top=6pt,bottom=6pt,
  before skip=9pt,after skip=9pt}
\newcommand{\eqE}{\mathrel{\equiv_E}}
\newcommand{\cl}[1]{[#1]_E}
\newcommand{\stern}{\mathbin{\ast}}
\newcommand{\probe}{\mathrel{\sim_D}}
\title{Was zweiseitige Produkte verraten\\[.6em]
  \large Ununterscheidbarkeit, Klassen und Quotienten in Halbgruppen\\[.4em]
  Experimentelle Lesefassung zu Band 28}
\author{Martin Kunze}
\date{7. Oktober 2026}
\begin{document}
\hypersetup{pageanchor=false}
\newgeometry{left=24mm,right=24mm,top=24mm,bottom=25mm}
\maketitle
\restoregeometry
\hypersetup{pageanchor=true,hypertexnames=false}
\chapter*{Was zweiseitige Produkte verraten}
\addcontentsline{toc}{chapter}{Was zweiseitige Produkte verraten}

Zwei Elemente einer Halbgruppe können verschieden sein und sich dennoch
in bestimmten Produkten genau gleich verhalten. Um dies zu untersuchen,
setzen wir ein Element zwischen zwei andere und vergleichen die Ergebnisse:
\[
  x\,\boxed{a}\,y\qquad\text{und}\qquad x\,\boxed{b}\,y.
\]
Die äußeren Faktoren bleiben gleich; nur der mittlere Faktor wird
ausgetauscht. Falls kein solches Produkt den Austausch bemerkt, wollen
wir die beiden mittleren Elemente als äquivalent ansehen.

Was bedeutet diese Ununterscheidbarkeit? Warum entstehen daraus Klassen?
Und können wir mit diesen Klassen wieder multiplizieren? Wir entwickeln
die Antworten an einer kleinen Multiplikationstafel. Danach folgen die
allgemeinen Beweise, weitere Beispiele und ein Vergleich mit den
Nulltests aus \emph{Was Nullprodukte verraten}. Ein abschließender
Ausblick ordnet die Konstruktion in die Halbgruppenforschung ein und
erläutert, was sich über ihren geschichtlichen Ursprung belegen lässt.

\Abschnitt{1. Welche Formel untersuchen wir?}{sec:formula}

Wir arbeiten mit einer nichtleeren Halbgruppe. Ihre Grundmenge nennen
wir \(C\), und ihr Produkt schreiben wir durch Nebeneinandersetzen.
\Rand{\(C\); \(ab\)}{Grundmenge der Halbgruppe; Produkt zweier Elemente.}
Das bedeutet: Zu \(a,b\in C\) gehört wieder ein Element \(ab\in C\),
und für alle \(a,b,c\in C\) gilt
\[
  (ab)c=a(bc).
\]
Diese Assoziativität erlaubt die Schreibweise \(abc\) ohne Angabe einer
Klammerung. Kommutativität setzen wir nicht voraus: Die Reihenfolge der
Faktoren bleibt entscheidend.

Die vorgeschlagene Relation lautet wörtlich
\Rand{\(E\); \(\equiv_E\)}{Parametermenge und Name der untersuchten Relation.}
\begin{equation}\label{eq:original}
 a\eqE b
 \quad:\Longleftrightarrow\quad
 \forall x,y\in C\ \forall e\in E:\ xay=xby,
 \qquad a,b\in C.
\end{equation}
Dabei kommt \(e\) in der Gleichung \(xay=xby\) nicht vor. Es bezeichnet
in dieser Formel insbesondere kein Einselement. Die Menge \(E\) wählt
auch keine äußeren Faktoren aus; diese werden aus \(C\) genommen.

Ist \(E\) nichtleer, können wir irgendein \(e\in E\) wählen. Dann
fordert die rechte Seite für jedes \(x,y\) die Gleichheit \(xay=xby\).
Umgekehrt gilt eine von \(e\) unabhängige wahre Gleichheit für jedes
\(e\in E\). Deshalb ist \eqref{eq:original} dann genau die Bedingung
\begin{equation}\label{eq:context}
  a\eqE b
  \quad\Longleftrightarrow\quad
  xay=xby\ \text{für alle }x,y\in C.
\end{equation}
Für verschiedene nichtleere Mengen \(E\) ergibt sich also dieselbe Relation.
Der Index bleibt hier erhalten, um an die Ausgangsnotation anzuschließen.

Ist dagegen \(E=\varnothing\), gibt es kein \(e\), für das etwas
geprüft werden müsste. Die Allquantifikation über \(E\) ist dann wahr,
unabhängig von \(a,b,x,y\). In diesem Fall sind alle Elemente von \(C\)
äquivalent. Wir behandeln diesen Randfall später ausdrücklich.
\textbf{Bis dahin sei \(E\ne\varnothing\).}

Die äußeren Faktoren \(x,y\) nennen wir einen \emph{zweiseitigen
Kontext}. Ein Kontext ist hier schlicht ein geordnetes Paar von
Halbgruppenelementen. Die Aussage \(a\eqE b\) bedeutet: Jeder dieser
Kontexte liefert für \(a\) und \(b\) denselben Produktwert.

\textbf{Warum betrachten wir gerade diese Beobachtungen?}
Manchmal interessiert uns ein Element vor allem durch seine Wirkung
zwischen zwei anderen Faktoren. Wir fragen dann: Kann es in dieser
mittleren Position durch ein anderes Element ersetzt werden, ohne dass
irgendein zugelassener Versuch ein anderes Ergebnis liefert? Genau
darauf antwortet unsere Relation.

Eine Klasse wird später alle Elemente mit dieser gemeinsamen Rolle
sammeln. Der Übergang zu Klassen vereinfacht die Beschreibung, indem
er Unterschiede zusammenfasst, die unsere zweiseitigen Tests nicht
erkennen. Wie das Multiplikationsverhalten dabei zu einer Operation
auf Klassen wird, müssen wir noch beweisen.

Diese Wahl von Beobachtungen gehört zur Fragestellung: Für einseitige
Produkte können Unterschiede wichtig bleiben, die in zweiseitigen
Produkten verschwinden. Wir behaupten daher nicht, dass äquivalente
Elemente in jeder denkbaren Untersuchung austauschbar sind. Das folgende
Beispiel zeigt sowohl den Nutzen als auch die Grenze unserer Tests.

\Abschnitt{2. Ein Beispiel mit fünf Elementen}{sec:example}

Unser durchgehendes Beispiel ist
\Rand{\(u,u^2,u^3,u^4,0\)}{Fünf verschiedene Elemente des Beispiels; \(0\) absorbiert jedes Produkt.}
\[
  C=\{u,u^2,u^3,u^4,0\}.
\]
Die Symbole bezeichnen fünf verschiedene Elemente. Die Potenzschreibweise
passt zu ihrer Multiplikation: Exponenten werden addiert, solange ihre
Summe kleiner als fünf ist; ab fünf ist das Ergebnis \(0\).
Außerdem gilt \(0a=a0=0\) für jedes \(a\in C\). Die vollständige Tafel ist
\[
\begin{array}{c|ccccc}
 \cdot&u&u^2&u^3&u^4&0\\\hline
 u   &u^2&u^3&u^4&0&0\\
 u^2 &u^3&u^4&0&0&0\\
 u^3 &u^4&0&0&0&0\\
 u^4 &0&0&0&0&0\\
 0   &0&0&0&0&0
\end{array}
\]
So ist \(u\cdot u=u^2\), \(u\cdot u^2=u^3\) und
\(u^2\cdot u^3=0\). Insbesondere hat \(0\) hier eine algebraische
Bedeutung: Es ist ein \emph{Nullelement}, das auf beiden Seiten absorbiert.

\textbf{Warum ist diese Operation assoziativ?}
Enthält ein Dreierprodukt \(0\), ergeben beide Klammerungen \(0\).
Betrachten wir sonst \(u^r,u^s,u^t\) mit \(r,s,t\in\{1,2,3,4\}\).
Für \(r+s+t<5\) bleiben auch beide Zwischensummen unter fünf; beide
Klammerungen liefern \(u^{r+s+t}\). Für \(r+s+t\geq5\) ist das Ergebnis
in beiden Klammerungen Null: Entweder ist schon das Zwischenprodukt
Null, oder das letzte Multiplizieren erreicht die Grenze fünf.
Damit sind alle Fälle erfasst.

\begin{Lesebeispiel}{Ein Kontext unterscheidet bereits drei Verhaltensweisen}
Wählen wir auf beiden Seiten \(u\). Beim Austausch des mittleren
Elements erhalten wir
\[
\begin{array}{c|ccccc}
 \text{mittleres Element }a&u&u^2&u^3&u^4&0\\\hline
 uau&u^3&u^4&0&0&0
\end{array}
\]
Der Kontext \((u,u)\) unterscheidet also \(u\) von \(u^2\) und beide
von den drei übrigen Elementen. Er unterscheidet aber \(u^3,u^4,0\)
nicht voneinander. Ob diese drei wirklich äquivalent sind, lässt sich
erst sagen, wenn \emph{alle} Kontexte verglichen worden sind.
\end{Lesebeispiel}

\Abschnitt{3. Was ein Test sieht und was er übersieht}{sec:tests}

Ein einzelner Test besteht darin, zwei äußere Faktoren \(x,y\) festzuhalten
und die Werte \(xay,xby\) zu vergleichen.
\Rand{\(a,b\); \((x,y)\)}{Kandidaten in der Mitte; geordnetes Paar der äußeren Testfaktoren.}
Zum Widerlegen von \(a\eqE b\) genügt \emph{ein} Kontext mit
verschiedenen Werten. Zum Beweisen der Äquivalenz müssen die Werte
in \emph{jedem} Kontext übereinstimmen.

Für das Fünferbeispiel können wir die zweite Aufgabe ohne 25 einzelne
Rechnungen erledigen. Steht \(u^3\) in der Mitte und sind die äußeren
Faktoren beide von Null verschieden, hat das Produkt einen Gesamtexponenten
von mindestens \(1+3+1=5\). Es ist deshalb Null. Für \(u^4\) beträgt
der Gesamtexponent mindestens \(1+4+1=6\), also ist das Produkt ebenfalls
Null. Mit einem äußeren Nullfaktor ist das
Ergebnis ohnehin Null. Folglich gilt
\[
  xu^3y=xu^4y=x0y=0\qquad(x,y\in C).
\]
Also sind \(u^3,u^4,0\) äquivalent. Der Test aus dem vorigen Abschnitt
zeigt zugleich, dass keines dieser Elemente zu \(u\) oder \(u^2\)
äquivalent ist und dass \(u\) und \(u^2\) verschiedenes Verhalten haben.
Die drei Gruppen sind damit vollständig bestimmt.

\textbf{Äquivalenz bedeutet hier nicht Gleichheit.}
Die Elemente \(u^3\) und \(0\) bleiben verschieden. Sogar in einem
Zweierprodukt lassen sie sich unterscheiden:
\[
  u\cdot u^3=u^4\ne0=u\cdot0.
\]
Unsere Tests schreiben aber einen Faktor auf \emph{beiden} Seiten vor.
Der zusätzliche rechte Faktor macht im Beispiel auch \(u^4\) zu Null.
Gerade diese Einschränkung lässt den Unterschied verschwinden.

Man darf deshalb nicht von \(a\eqE b\) auf \(xa=xb\) oder \(ay=by\)
schließen. Ebenso wenig darf man einen äußeren Faktor einfach weglassen.
Ein \enquote{leerer Kontext} steht in der Definition nicht zur Verfügung.

\Abschnitt{4. Warum dies eine Äquivalenzrelation ist}{sec:equivalence}

Eine Äquivalenzrelation muss reflexiv, symmetrisch und transitiv sein.
Wir überprüfen jede dieser Eigenschaften direkt. Dabei bleiben die
äußeren Faktoren in jedem Vergleich dieselben.

\textbf{Reflexivität.}
Sei \(a\in C\). Für beliebige \(x,y\in C\) gilt
\[
  xay=xay.
\]
Die Bedingung \eqref{eq:context} ist also erfüllt: \(a\eqE a\).
Jedes Element hat dasselbe Verhalten wie es selbst.

\textbf{Symmetrie.}
Seien \(a,b\in C\) und \(a\eqE b\). Für beliebige \(x,y\in C\)
wissen wir dann \(xay=xby\). Gleichheit lässt sich umkehren, also gilt
\(xby=xay\). Da dies für jedes Paar \(x,y\) gilt, folgt \(b\eqE a\).

\textbf{Transitivität.}
Seien \(a,b,c\in C\) mit \(a\eqE b\) und \(b\eqE c\). Wir wählen
wieder beliebige \(x,y\in C\). Die erste Voraussetzung liefert die
erste Gleichheit, die zweite die zweite:
\[
  xay=xby=xcy.
\]
Daraus folgt \(xay=xcy\). Die äußeren Faktoren waren beliebig; also
gilt \(a\eqE c\).

Damit ist \(\equiv_E\) eine Äquivalenzrelation. Dass zwei Elemente
dieselben Antworten geben, kann also tatsächlich als eine gemeinsame
Verhaltensart aufgefasst werden.

\textbf{Wo wurde Assoziativität benutzt?}
In diesen drei Schlüssen überhaupt nicht. Es genügt, jedem mittleren
Element und jedem Testpaar einen eindeutig bestimmten Wert zuzuordnen.
Würde man in einer beliebigen Menge mit binärer Operation ausdrücklich
\((xa)y\) vergleichen, ergäbe auch die Gleichheit dieser Werte eine
Äquivalenzrelation. Assoziativität brauchen wir später, um aus
der Äquivalenz eine verträgliche Multiplikation der Klassen zu gewinnen.

\Abschnitt{5. Das ganze Verhalten als Profil}{sec:profiles}

Statt die einzelnen Gleichungen immer wieder auszuschreiben, können
wir alle Testantworten zu einer Funktion zusammenfassen. Für \(a\in C\)
setzen wir
\Rand{\(C\times C\); \(P_a\)}{Menge der geordneten Testpaare; Funktion aller Antworten für \(a\).}
\[
  P_a\colon C\times C\longrightarrow C,
  \qquad P_a(x,y)=xay.
\]
Wir nennen \(P_a\) das \emph{Kontextprofil} von \(a\). Eine Funktion
ist hier durch sämtliche Werte zu ihren Argumenten bestimmt. Folglich
ist unsere Relation genau die Gleichheit der Profile:
\[
  a\eqE b\quad\Longleftrightarrow\quad P_a=P_b.
\]

Ein Profil hält nicht nur fest, \emph{welche} Werte irgendwann auftreten.
Es hält auch fest, \emph{welcher Kontext} welchen Wert liefert. Zwei
Profile müssen an jedem einzelnen Testpaar übereinstimmen. Bloß dieselbe
Menge möglicher Ausgaben zu haben, wäre eine schwächere Bedingung.

Im endlichen Fall kann man ein Profil als Tabelle lesen: Die Zeile
bestimmt \(x\), die Spalte \(y\), und im Feld steht \(xay\). Für das
Fünferbeispiel sind alle nicht aufgeführten Felder Null; der möglicherweise
von Null verschiedene Teil liegt bei \(x,y\in\{u,u^2\}\):
\[
\begin{array}{c|cc}
 P_u&u&u^2\\\hline
 u&u^3&u^4\\
 u^2&u^4&0
\end{array}
\qquad
\begin{array}{c|cc}
 P_{u^2}&u&u^2\\\hline
 u&u^4&0\\
 u^2&0&0
\end{array}
\]
Die drei Profile \(P_{u^3},P_{u^4},P_0\) sind dagegen überall Null.
Ein äußerer Faktor \(u^3,u^4\) oder \(0\) macht in diesem Beispiel
jedes Dreierprodukt Null. Damit sind auch die weggelassenen Zeilen und
Spalten beschrieben.

\Abschnitt{6. Was eine Äquivalenzklasse ist}{sec:classes}

Für \(a\in C\) definieren wir
\Rand{\([a]_E\)}{Klasse von \(a\): alle Grundelemente mit demselben Kontextprofil.}
\[
  \cl a=\{b\in C\mid b\eqE a\}.
\]
Die Klasse ist eine \emph{Teilmenge} von \(C\), während \(a\) ein
\emph{Element} von \(C\) ist. Wegen der Reflexivität gehört \(a\)
selbst zu seiner Klasse. Die Klasse ist daher nie leer.

\textbf{Warum dürfen verschiedene Vertreter dieselbe Klasse benennen?}
Nehmen wir \(a\eqE b\) an. Für \(c\in\cl a\) gilt \(c\eqE a\),
also mit \(a\eqE b\) auch \(c\eqE b\). Daher gehört \(c\) zu
\(\cl b\). Das zeigt \(\cl a\subseteq\cl b\). Umgekehrt ergibt
\(c\eqE b\) zusammen mit \(b\eqE a\) die Zugehörigkeit
\(c\in\cl a\). Somit gilt \(\cl a=\cl b\).

Gilt umgekehrt \(\cl a=\cl b\), dann gehört \(a\) wegen der
Reflexivität zu \(\cl a\) und damit zu \(\cl b\). Nach Definition
heißt das \(a\eqE b\). Wir haben also beide Richtungen bewiesen:
\[
  a\eqE b\quad\Longleftrightarrow\quad\cl a=\cl b.
\]

\textbf{Warum überlappen verschiedene Klassen nicht?}
Angenommen, \(c\) gehört sowohl zu \(\cl a\) als auch zu \(\cl b\).
Dann gilt \(c\eqE a\) und \(c\eqE b\). Durch Symmetrie und
Transitivität folgt \(a\eqE b\), also \(\cl a=\cl b\).
Zwei Klassen sind folglich entweder gleich oder disjunkt. Zusammen
überdecken sie \(C\), denn jedes Element gehört zu seiner eigenen Klasse.

\begin{Lesebeispiel}{Fünf Grundelemente, drei Klassen}
Im Fünferbeispiel sind die Klassen genau
\[
 \cl u=\{u\},\qquad
 \cl{u^2}=\{u^2\},\qquad
 \cl0=\{0,u^3,u^4\}.
\]
Dabei gilt \(\cl{u^3}=\cl{u^4}=\cl0\). Die letzten drei Namen
bezeichnen dieselbe Klasse. Die Klasse \(\cl0\) enthält drei Elemente;
bei der Zählung der \emph{Klassen} wird sie aber nur einmal gezählt.
\end{Lesebeispiel}

Die Klassen selbst müssen nicht unter der ursprünglichen Multiplikation
abgeschlossen sein. Im Beispiel gehört \(u\) zur Klasse \(\{u\}\),
aber \(u\cdot u=u^2\) gehört nicht zu ihr. Eine Klasse ist deshalb
im Allgemeinen keine Unterhalbgruppe. Erst die \emph{Menge aller
Klassen} wird später mit einer eigenen Operation zur Halbgruppe.

\Abschnitt{7. Können wir Klassen multiplizieren?}{sec:congruence}

Wir wollen ein Produkt zweier Klassen dadurch bestimmen, dass wir
Vertreter wählen, sie in \(C\) multiplizieren und danach die Klasse
des Ergebnisses nehmen. Das kann nur funktionieren, wenn die Antwort
von der Wahl der Vertreter unabhängig ist.

Sei zunächst \(a\eqE b\), und sei \(c\in C\). Wir beweisen zuerst
\(ca\eqE cb\). Für beliebige \(x,y\in C\) gilt durch Assoziativität
\[
  x(ca)y=(xc)ay=(xc)by=x(cb)y.
\]
Die mittlere Gleichheit ist gerade die Voraussetzung \(a\eqE b\),
angewandt mit dem linken Testfaktor \(xc\) und dem rechten Testfaktor
\(y\). Das ist erlaubt, weil \(xc\in C\).

Ebenso zeigen wir \(ac\eqE bc\). Diesmal rechnen wir
\[
  x(ac)y=xa(cy)=xb(cy)=x(bc)y.
\]
Hier benutzen wir \(a\eqE b\) mit den Testfaktoren \(x\) und
\(cy\). Auch \(cy\) ist ein Element von \(C\).

Eine Äquivalenzrelation mit diesen beiden Verträglichkeiten heißt
\emph{Halbgruppenkongruenz}. Unsere Relation ist also sogar eine
Kongruenz. Sie bleibt erhalten, wenn wir links oder rechts mit einem
festen Halbgruppenelement multiplizieren.

\Needspace{8\baselineskip}
\textbf{Und wenn sich beide Vertreter ändern?}
Seien \(a\eqE a'\) und \(b\eqE b'\). Rechtsmultiplikation mit \(b\)
liefert \(ab\eqE a'b\). Linksmultiplikation mit \(a'\) liefert
\(a'b\eqE a'b'\). Mit Transitivität folgt
\begin{equation}\label{eq:congruence}
  ab\eqE a'b'.
\end{equation}
Genau diese Aussage brauchen wir für das Produkt von Klassen. Hier
arbeiten Assoziativität und Abgeschlossenheit der Halbgruppe zusammen:
Ein benachbarter Faktor kann in den äußeren Testfaktor aufgenommen werden.

Das erklärt auch den Austausch in längeren Produkten: Solange mindestens
ein Faktor links und mindestens ein Faktor rechts vom auszutauschenden
Element steht, können wir jeden äußeren Block zu einem Element von
\(C\) zusammenfassen. Dann gilt unmittelbar der Kontextvergleich.
Fehlt ein solcher Block, steht uns dieser Schluss über genaue
Produktwerte nicht zur Verfügung.

\Abschnitt{8. Die Quotientenhalbgruppe}{sec:quotient}

Die Menge der verschiedenen Klassen schreiben wir als
\Rand{\(Q=C/{\equiv_E}\)}{Menge aller verschiedenen Äquivalenzklassen.}
\[
  Q=C/{\equiv_E}=\{\cl a\mid a\in C\}.
\]
Mehrfach auftretende Klassennamen ergeben nur ein Element dieser Menge.
Auf \(Q\) definieren wir die neue Operation
\Rand{\(\ast\)}{Produkt in \(Q\): Vertreter multiplizieren und die Ergebnisklasse nehmen.}
\[
  \cl a\stern\cl b:=\cl{ab}.
\]
Wählt man \(a',b'\) aus denselben Klassen, so gilt
\(a\eqE a'\) und \(b\eqE b'\). Nach \eqref{eq:congruence} ist
\(ab\eqE a'b'\), also \(\cl{ab}=\cl{a'b'}\). Die Definition
liefert deshalb unabhängig von den Vertretern dieselbe Klasse.
Damit ist \(\ast\) eine wohldefinierte binäre Operation auf \(Q\).

Auch ihre Assoziativität können wir direkt nachrechnen:
\[
\begin{aligned}
 (\cl a\stern\cl b)\stern\cl c
    &=\cl{(ab)c},\\
 \cl a\stern(\cl b\stern\cl c)
    &=\cl{a(bc)}.
\end{aligned}
\]
Da \((ab)c=a(bc)\) in \(C\), sind die rechten Seiten gleich. Somit
ist \((Q,\ast)\) eine Halbgruppe: die \emph{Quotientenhalbgruppe}.

\textbf{Klassenprodukt und Mengenprodukt sind verschiedene Operationen.}
Die Klassen sind Teilmengen von \(C\). Man könnte daher auch alle
Einzelprodukte aus ihnen sammeln:
\Rand{\(AB\)}{Mengenprodukt: Menge aller \(ab\) mit \(a\in A\), \(b\in B\).}
\[
  AB:=\{ab\mid a\in A,\ b\in B\}\qquad(A,B\subseteq C).
\]
Dieses Mengenprodukt muss keine ganze Klasse ergeben. Es liegt aber
für \(A=\cl a\) und \(B=\cl b\) vollständig in \(\cl{ab}\).
Denn jeder seiner Werte entsteht aus Vertretern derselben beiden Klassen.

Im Fünferbeispiel gilt beispielsweise
\[
  \cl0\stern\cl0=\cl0=\{0,u^3,u^4\},
  \qquad \cl0\,\cl0=\{0\}.
\]
Die zweite Gleichung benutzt das Mengenprodukt; jedes Paar aus
\(\{0,u^3,u^4\}\) hat das Produkt Null. Die Ergebnisklasse enthält
also mehr Elemente als das tatsächlich gebildete Mengenprodukt.

\Abschnitt{9. Die neue Tafel unseres Beispiels}{sec:quotientexample}

Für die drei Klassen verwenden wir nun kürzere Namen:
\Rand{\(K_1,K_2,K_0\)}{Die Klassen von \(u\), \(u^2\) und \(0\) im Fünferbeispiel.}
\[
  K_1=\cl u,\qquad K_2=\cl{u^2},\qquad K_0=\cl0.
\]
Der Index ist ein Name; insbesondere enthält \(K_0\) drei
Grundelemente. Aus der ursprünglichen Multiplikation erhalten wir
\[
\begin{array}{c|ccc}
 \ast&K_1&K_2&K_0\\\hline
 K_1&K_2&K_0&K_0\\
 K_2&K_0&K_0&K_0\\
 K_0&K_0&K_0&K_0
\end{array}
\]
So ist \(K_1\ast K_1=\cl{u^2}=K_2\), während
\(K_1\ast K_2=\cl{u^3}=K_0\). Die Klasse \(K_0\) ist ein
Nullelement der neuen Halbgruppe: Ihr Produkt mit jeder Klasse ist
auf beiden Seiten wieder \(K_0\).

\textbf{Eine konkrete Änderung des Vertreters.}
Um \(K_1\ast K_0\) zu berechnen, können wir in \(K_0\) das Element
\(0\) oder das Element \(u^3\) wählen. Im ersten Fall entsteht
\(u\cdot0=0\), im zweiten \(u\cdot u^3=u^4\). Die Grundelemente
\(0,u^4\) sind verschieden, aber ihre Klassen sind gleich. Beide
Rechnungen ergeben daher \(K_0\). Vertreterunabhängigkeit fordert
die Gleichheit der Ergebnisklassen, nicht die Gleichheit der Einzelprodukte.

\textbf{Was wurde zusammengefasst?}
Aus fünf Grundelementen sind drei Elemente von \(Q\) geworden.
Die neue Tafel beschreibt weiterhin jedes ursprüngliche Produkt bis
auf seine Klasse. Sie enthält aber weniger Information: Ob ein
Produkt ursprünglich \(0\), \(u^3\) oder \(u^4\) war, ist an
der Klasse \(K_0\) allein nicht mehr abzulesen.

\Abschnitt{10. Die Abbildung von Elementen zu ihren Klassen}{sec:projection}

Die Zuordnung zur eigenen Klasse ist eine Funktion
\Rand{\(\pi\)}{Kanonische Projektion: ordnet jedem Grundelement seine Klasse zu.}
\[
  \pi\colon C\longrightarrow Q,\qquad \pi(a)=\cl a.
\]
Sie ist surjektiv, denn jede Klasse in \(Q\) hat nach Definition
mindestens einen Vertreter. Außerdem erhält sie Produkte:
\[
  \pi(ab)=\cl{ab}=\cl a\stern\cl b=\pi(a)\stern\pi(b).
\]
Eine Funktion mit dieser Produkteigenschaft heißt
\emph{Halbgruppenhomomorphismus}. Die Projektion ist also ein
surjektiver Homomorphismus.

Die Elemente, die \(\pi\) auf denselben Wert abbildet, sind genau
die äquivalenten Elemente:
\[
  \pi(a)=\pi(b)\quad\Longleftrightarrow\quad a\eqE b.
\]
Im Fünferbeispiel ordnet \(\pi\) sowohl \(0\) als auch \(u^3,u^4\)
der Klasse \(K_0\) zu. Die Funktion ist dort also nicht injektiv.
Allgemein ist sie genau dann injektiv, wenn jede Klasse eine Einermenge ist.

Daneben können wir jedem Element sein Profil zuordnen. Dazu schreiben wir
\Rand{\(\mathcal P\); \(\rho\)}{Vorkommende Profile; Zuordnung \(a\mapsto P_a\).}
\[
  \mathcal P=\{P_a\mid a\in C\},\qquad
  \rho\colon C\to\mathcal P,\quad\rho(a)=P_a.
\]
Auch \(\rho\) fasst genau die äquivalenten Elemente zusammen.
Deshalb entspricht jeder Klasse genau ein vorkommendes Profil:
\[
  \cl a\longmapsto P_a.
\]
Diese Zuordnung ist wohldefiniert, weil äquivalente Vertreter dasselbe
Profil haben. Sie ist injektiv, weil gleiche Profile äquivalente
Vertreter und damit gleiche Klassen ergeben. Sie ist surjektiv,
weil jedes Profil in \(\mathcal P\) das Profil irgendeines Elements ist.
Klassen und vorkommende Profile sind somit zwei Beschreibungen derselben
Aufteilung. Ein Produkt von Profilen wurde dadurch noch nicht definiert;
das Produkt von Klassen ist bereits durch \(\ast\) festgelegt.

\Abschnitt{11. Zwei äußerste Fälle}{sec:extremes}

Wie stark die Relation zusammenfasst, hängt von der Halbgruppe ab.
Sie kann gar keine verschiedenen Elemente identifizieren oder alle.

\textbf{Mit einem Einselement erkennt ein einziger Kontext jedes Element.}
Hat \(C\) ein Einselement \(1\), gilt \(1a=a1=a\) für jedes \(a\in C\).
\Rand{\(1\)}{Einselement: lässt beim Multiplizieren jedes Element unverändert.}
Aus \(a\eqE b\) folgt durch den Test \((1,1)\)
\[
  a=1a1=1b1=b.
\]
Umgekehrt gilt aus \(a=b\) selbstverständlich \(a\eqE b\).
In einem Monoid, also einer Halbgruppe mit Einselement, ist unsere
Relation daher genau die Gleichheitsrelation. Alle Klassen sind Einermengen.
Hier wird weiterhin \(E\ne\varnothing\) vorausgesetzt.

Auch eine nichtleere Halbgruppe, in der auf beiden Seiten gekürzt werden
darf, hat nur Einermengen als Klassen. Wählen wir irgendwelche \(x,y\in C\).
Aus \(xay=xby\) folgt durch Kürzen des linken Faktors \(x\) zunächst
\(ay=by\); durch Kürzen des rechten Faktors \(y\) folgt \(a=b\).
Ein Einselement ist für dieses Argument nicht nötig.

\textbf{Wenn alle Zweierprodukte Null sind, erkennt kein Kontext einen Unterschied.}
In einer Nullhalbgruppe gilt \(ab=0\) für jedes Paar \(a,b\).
Dann ist auch jedes Dreierprodukt Null. Für alle \(a,b,x,y\) gilt also
\(xay=0=xby\). Es gibt nur eine Klasse, nämlich \(C\) selbst.
Der Quotient hat ein einziges Element, und dessen Produkt mit sich selbst
ist wieder dieses Element.

Eine solche Klasse sagt nichts darüber aus, wie viele Grundelemente
\(C\) hatte: Nullhalbgruppen mit zwei, fünf oder beliebig vielen Elementen
haben unter dieser Relation jeweils nur eine Klasse.

\Abschnitt{12. Eine einzige Klasse ohne Nullelement}{sec:rectangular}

Dass alle Elemente äquivalent sind, muss nicht daran liegen, dass alle
Produkte Null werden. Ein zweites Beispiel lässt sich als Rechteck lesen.
\Rand{\(I,J\); \(B=I\times J\)}{Zeilen- und Spaltenmenge; Grundmenge des Rechteckbeispiels.}
Seien \(I,J\) nichtleere Mengen. Auf \(B=I\times J\) setzen wir
\[
  (i,j)(k,\ell)=(i,\ell).
\]
Das Ergebnis übernimmt die Zeile des linken und die Spalte des rechten
Faktors. Für drei Faktoren ergibt jede Klammerung
\[
  (i,j)(k,\ell)(m,n)=(i,n),
\]
also ist die Operation assoziativ. Eine solche Halbgruppe heißt
\emph{Rechtecksband}.

Nehmen wir \(x=(i,j)\), einen beliebigen mittleren Faktor
\(a=(k,\ell)\) und \(y=(m,n)\). Dann ist
\[
  xay=(i,n).
\]
Die Zeile und Spalte des mittleren Faktors werden beide überschrieben.
Der Wert hängt von \(x,y\) ab, aber überhaupt nicht von \(a\).
Somit sind alle Elemente äquivalent, obwohl das gemeinsame Profil
im Allgemeinen nicht konstant ist.

\begin{Lesebeispiel}{Vier Felder, ein Kontextprofil}
Für \(I=J=\{1,2\}\) hat \(B\) die vier Elemente
\((1,1),(1,2),(2,1),(2,2)\). Mit linkem Faktor \((1,2)\) und rechtem
Faktor \((2,1)\) ergibt jeder mittlere Faktor den Wert \((1,1)\).
Mit linkem Faktor \((2,1)\) und rechtem Faktor \((1,2)\) ergibt
jeder mittlere Faktor dagegen \((2,2)\). Die Kontexte liefern also
unterschiedliche Werte, behandeln aber alle mittleren Elemente gleich.

Dieses Beispiel hat keine Null. Wäre \(z=(i,j)\) eine Null, müsste
\(zx=z\) für jedes \(x\) gelten. Aber \(z(k,1)=(i,1)\) und
\(z(k,2)=(i,2)\) sind verschieden. Schon die rechte Absorption ist unmöglich.
\end{Lesebeispiel}

Eine einzige Klasse bedeutet deshalb: \emph{Jeder feste Kontext ist
vom mittleren Element unabhängig.} Sie bedeutet nicht: \emph{Alle
Kontexte haben denselben Wert.}

\Abschnitt{13. Der Vergleich mit Nullproduktprofilen}{sec:nullprofiles}

In \emph{Was Nullprodukte verraten} werden andere Fragen gestellt.
Für einzelne Elemente in einer Halbgruppe mit Null lauten sie:
Ist \(ax=0\)? Und ist \(xa=0\)? Dort werden auch nichtleere
Teilmengen als Kandidaten und Testmengen betrachtet. Hier genügt
für den Vergleich zunächst die Elementebene.

Unser Kontexttest fragt dagegen nach dem \emph{genauen Wert} von
\(xay\) und verlangt einen Faktor auf jeder Seite. Das sind andere
Beobachtungen. Dass unser Test einen genauen Wert statt nur Ja oder
Nein liefert, macht ihn nicht automatisch unterscheidungskräftiger:
Der zweite äußere Faktor kann zuvor sichtbare Unterschiede auslöschen.

\Needspace{22\baselineskip}
Für das folgende Dreierbeispiel nennen wir die Grundmenge \(N\).
\Rand{\(N\)}{Dreielementige Halbgruppe des Vergleichs mit Nullproduktprofilen.}

\begin{Lesebeispiel}{Dasselbe Dreierbeispiel, zwei verschiedene Aufteilungen}
Sei \(N=\{0,c,a\}\) mit \(a^2=c\), alle anderen Zweierprodukte
seien Null. Dies ist das Dreierbeispiel der Nullprodukt-Lesefassung.
\[
\begin{array}{c|ccc}
 \cdot&0&c&a\\\hline
 0&0&0&0\\
 c&0&0&0\\
 a&0&0&c
\end{array}
\]
Jedes Zweierprodukt liegt in \(\{0,c\}\), und beide Elemente ergeben
mit jedem weiteren Faktor Null. Daher sind alle Dreierprodukte Null;
das begründet zugleich die Assoziativität.

Für unsere Kontextrelation gibt es auf \(N\) nur eine Klasse:
\(x0y=xcy=xay=0\) für jedes \(x,y\in N\). Beim einseitigen
Nulltest mit \(a\) sieht man hingegen einen Unterschied:
\[
  0a=ca=0,\qquad aa=c\ne0.
\]
Auf der Ebene der Grundelemente haben \(0,c\) dasselbe Nulltestverhalten,
während \(a\) ein anderes hat. Auf der Ebene nichtleerer Teilmengen
entstehen in der Vorlage entsprechend die zwei Familien: Mengen
ohne \(a\) und Mengen mit \(a\).
\end{Lesebeispiel}

Die beiden Lesefassungen teilen also eine Vorgehensweise: Man legt
Beobachtungen fest, vergleicht alle Antworten und fasst gleiches
Verhalten zu Klassen zusammen. Welche Klassen entstehen, hängt aber
von der genauen Wahl der Beobachtungen ab. Hier vergleichen wir
zweiseitige Produktwerte, dort linke und rechte Nullprodukte.

\Abschnitt{14. Was passiert mit einem hinzugefügten Einselement?}{sec:identity}

In der Halbgruppentheorie lässt man in anderen Zusammenhängen auch
äußere Faktoren aus einer Halbgruppe mit hinzugefügter Eins zu.
\Rand{\(C^1\)}{\(C\) mit Einselement; eine neue Eins wird nur bei Bedarf angefügt.}
Wir schreiben \(C^1\) für \(C\), falls bereits eine Eins existiert,
und sonst für \(C\) mit einem neuen Element \(1\), das
\(1a=a1=a\) erfüllt. Die bisherige Multiplikation auf \(C\) bleibt erhalten.

Vergleicht man nun
\[
  xay=xby\qquad\text{für alle }x,y\in C^1,
\]
so darf man \(x=y=1\) setzen. Die Bedingung erzwingt \(a=b\).
Für \(a,b\in C\) ist diese veränderte Relation deshalb immer die
Gleichheitsrelation, auch wenn die ursprüngliche Halbgruppe keine Eins hatte.

Im Fünferbeispiel sind \(u^3\) und \(0\) für Kontexte aus \(C\)
äquivalent. Mit der neuen Eins sind sie sofort zu unterscheiden:
\[
  1u^3 1=u^3\ne0=1\cdot0\cdot1.
\]
Die Wahl der zulässigen äußeren Faktoren gehört somit wesentlich
zur Definition. Die Formulierung \enquote{in allen Kontexten} muss
immer sagen, welche Kontexte zugelassen sind.

Für die hier untersuchte Relation bleiben die Faktoren in \(C\).
Wir führen keinen leeren linken oder rechten Faktor ein. Die
Assoziativität erlaubt das Umklammern vorhandener Faktoren; sie
erzeugt kein Einselement und erlaubt nicht, Faktoren zu entfernen.

\Abschnitt{15. Wenn nur ausgewählte Kontexte erlaubt sind}{sec:restricted}

Bisher ist \(C\) die gesamte Grundmenge der Halbgruppe. Man kann
stattdessen eine Teilmenge \(D\subseteq C\) als Vorrat äußerer
Testfaktoren festlegen und eine neue Relation definieren:
\Rand{\(D\); \(\sim_D\)}{Ausgewählte Testfaktoren; Gleichheit aller Tests mit Faktoren aus \(D\).}
\[
  a\probe b
  \quad:\Longleftrightarrow\quad
  xay=xby\ \text{für alle }x,y\in D.
\]
Reflexivität, Symmetrie und Transitivität gelten mit denselben
Argumenten wie zuvor. Eine beliebige Auswahl von Tests erzeugt also
weiterhin eine Äquivalenzrelation.

Für die Kongruenzeigenschaft reicht dieses Argument nicht aus.
In der Rechnung \(x(ca)y=(xc)ay\) müsste \(xc\) wieder ein
\emph{zugelassener} Testfaktor sein. Aus \(x\in D\) und \(c\in C\)
folgt das für eine beliebige Teilmenge \(D\) nicht.

Eine hinreichende Bedingung ist
\[
  xc\in D\quad\text{und}\quad cx\in D
  \qquad(x\in D,\ c\in C).
\]
Eine nichtleere Teilmenge mit dieser Eigenschaft heißt
\emph{zweiseitiges Ideal}. Unter dieser Bedingung funktionieren die
beiden Kongruenzrechnungen aus Abschnitt~7 auch für \(\sim_D\).
Für \(D=\varnothing\) ist die Relation ohnehin universell und ebenfalls
eine Kongruenz.

\textbf{Ein Gegenbeispiel ohne diese Abgeschlossenheit.}
Betrachten wir die fünf Matrizen
\Rand{\(M\); \(e_{ij}\)}{Nullmatrix und vier Matrixeinheiten; \(e_{ij}\) hat genau an Stelle \((i,j)\) eine Eins.}
\[
  M=\{0,e_{11},e_{12},e_{21},e_{22}\}
\]
unter gewöhnlicher Matrixmultiplikation. Die Matrix \(e_{ij}\) hat
an Stelle \((i,j)\) den Eintrag Eins und sonst nur Nullen. Ihre Produkte
sind \(e_{ij}e_{k\ell}=e_{i\ell}\), wenn \(j=k\), und sonst Null.
Die Menge ist unter der assoziativen Matrixmultiplikation abgeschlossen.
Konkret sind die vier von Null verschiedenen Matrizen
\[
\begin{aligned}
 e_{11}&=\begin{pmatrix}1&0\\0&0\end{pmatrix},&
 e_{12}&=\begin{pmatrix}0&1\\0&0\end{pmatrix},\\[.4em]
 e_{21}&=\begin{pmatrix}0&0\\1&0\end{pmatrix},&
 e_{22}&=\begin{pmatrix}0&0\\0&1\end{pmatrix}.
\end{aligned}
\]

Wählen wir nur \(D=\{e_{11}\}\) als Testvorrat. Dann gilt
\[
  e_{11}e_{12}e_{11}=0=e_{11}0e_{11},
\]
also \(e_{12}\probe0\). Rechtsmultiplikation mit \(e_{21}\) ergibt aber
\[
  e_{12}e_{21}=e_{11},\qquad 0e_{21}=0.
\]
Diese Ergebnisse sind nicht äquivalent für unseren Test, denn
\[
  e_{11}e_{11}e_{11}=e_{11}\ne0=e_{11}0e_{11}.
\]
Somit ist \(\sim_D\) hier keine Kongruenz. Klassen können also
existieren, ohne dass ihre Vertreterprodukte eine wohldefinierte
Quotientenmultiplikation ergeben.

\Abschnitt{16. Warum ein zweiter Quotient noch kleiner sein kann}{sec:iterate}

Die Klassen unterscheiden sich durch ihre Profile in der ursprünglichen
Halbgruppe. Daraus folgt aber nicht, dass sie sich auch durch
zweiseitige Tests \emph{innerhalb des Quotienten} unterscheiden lassen.
Die Ausgabewerte der Tests werden dort selbst zu Klassen zusammengefasst.

Im ursprünglichen Fünferbeispiel unterschied der Kontext \((u,u)\)
die mittleren Elemente \(u,u^2,0\): Er lieferte \(u^3,u^4,0\).
Diese drei Ausgaben liegen jedoch alle in der Klasse \(K_0\).
Nach dem Übergang zur Quotientenhalbgruppe werden die drei Testantworten
deshalb gleich:
\[
  K_1\stern K_1\stern K_1
  =K_1\stern K_2\stern K_1
  =K_1\stern K_0\stern K_1=K_0.
\]
Mehr noch: Jedes Dreierprodukt im Quotienten ist \(K_0\).
Nach seiner Tafel ist \(K_1\ast K_1=K_2\) das einzige Zweierprodukt,
das nicht \(K_0\) ist. Der weitere Faktor macht auch daraus \(K_0\),
auf beiden Seiten. Also liefern alle mittleren Quotientenelemente
in jedem Quotientenkontext denselben Wert.

Wenden wir die gleiche Konstruktion erneut an, entstehen daher
aus den drei Elementen des ersten Quotienten nur noch ein Element:
\[
  \underbrace{5\ \text{Grundelemente}}_{C}
  \ \longrightarrow\
  \underbrace{3\ \text{Klassen}}_{Q}
  \ \longrightarrow\
  1\ \text{Klasse im zweiten Quotienten}.
\]
Der erste Schritt war korrekt und vollständig. Der zweite Schritt
benutzt aber neue Tests mit neuen Ausgabewerten. Die Gleichheit zweier
Ausgaben in \(C\) ist eine strengere Forderung als die Gleichheit
ihrer Klassen in \(Q\).

Bei einer endlichen Halbgruppe kann man diese Konstruktion wiederholen.
Sobald sie verschiedene Elemente identifiziert, sinkt die Elementanzahl
streng. Das kann nur endlich oft geschehen. Bleibt die Anzahl gleich,
sind alle Klassen Einermengen; die Projektion ist dann ein Isomorphismus,
und ein weiterer Schritt ändert die Struktur nicht mehr. Ein stabiles
Ergebnis ist also bei endlichen Halbgruppen erreichbar, muss aber
nicht schon der erste Quotient sein.

\Abschnitt{17. Der leere Parameter und die vollständige Aussage}{sec:empty}

Zum Schluss kehren wir zur wörtlichen Formel \eqref{eq:original} zurück.
Für \(E\ne\varnothing\) haben wir bewiesen, dass sie die Gleichheit
aller zweiseitigen Kontextprofile beschreibt und eine Kongruenz ist.

Für \(E=\varnothing\) gilt \(a\eqE b\) dagegen für alle \(a,b\in C\).
Reflexivität, Symmetrie und Transitivität sind unmittelbar erfüllt.
Auch die Kongruenzbedingung gilt: Jedes beliebige Paar von Produkten
liegt wieder in \(C\) und ist deshalb äquivalent. Es gibt genau eine
Klasse, nämlich \(C\), und die Quotientenhalbgruppe hat ein Element.

Die Klasse wird in diesem Fall nicht durch gleiche Kontextprofile
bestimmt. Beispielsweise können in einem Monoid zwei verschiedene
Elemente durch \((1,1)\) unterschieden werden und trotzdem gemäß der
wörtlichen Formel für leeres \(E\) äquivalent sein. Die Aussagen
\(a\eqE b\Longleftrightarrow P_a=P_b\) und
\enquote{im Monoid sind alle Klassen Einermengen} brauchen daher
die ausdrücklich genannte Voraussetzung \(E\ne\varnothing\).

\begin{Lesebeispiel}{Die Aussage für beliebiges \(E\)}
Sei \(C\) eine nichtleere Halbgruppe, \(E\) eine beliebige Menge,
und die Relation sei durch \eqref{eq:original} definiert.
Dann ist \(\equiv_E\) eine Halbgruppenkongruenz. Ihre Klassen bilden
mit \(\cl a\ast\cl b=\cl{ab}\) eine Halbgruppe, und
\(a\mapsto\cl a\) ist ein surjektiver Homomorphismus.

Für nichtleeres \(E\) sind die Klassen genau die Gruppen gleicher
Kontextprofile. Für leeres \(E\) gibt es unabhängig von der
Multiplikation nur die eine Klasse \(C\).
\end{Lesebeispiel}

In dieser Fassung wurde \(C\) als gesamte Halbgruppen-Grundmenge
verwendet. Falls \(C\) stattdessen nur einen Testvorrat innerhalb einer
größeren Halbgruppe bezeichnet, gelten die Einschränkungen aus
Abschnitt~15. Falls der Buchstabe \(e\) in einer anderen beabsichtigten
Formel tatsächlich in einem Produkt vorkommt, muss diese andere
Formel gesondert untersucht werden.

\Abschnitt{18. Kleine Experimente zum Weiterdenken}{sec:exercises}

Die folgenden Aufgaben lassen sich mit den bisherigen Rechnungen
bearbeiten. In diesen Aufgaben sei wieder \(E\ne\varnothing\).
Kurze Lösungen stehen unmittelbar darunter.

\textbf{1. Eine frühere Nullgrenze.}
Nehmen wir \(\{v,v^2,v^3,0\}\) mit derselben Potenzregel, diesmal
\Rand{\(v\)}{Erzeugendes Element des Viererbeispiels in Aufgabe 1.}
aber mit \(v^4=0\). Welche Kontextklassen entstehen?

\textbf{Lösung.}
\(v\) wird vom Kontext \((v,v)\) auf \(v^3\ne0\) abgebildet.
Bei \(v^2,v^3,0\) ist jedes Dreierprodukt Null. Die Klassen sind
also \(\{v\}\) und \(\{v^2,v^3,0\}\). Im Quotienten ist bereits
jedes Zweierprodukt die zweite Klasse, weil \(v\cdot v=v^2\)
zu ihr gehört. Ein erneuter Kontextquotient hat nur ein Element.

\textbf{2. Ein Test zur Widerlegung.}
Warum genügt im Fünferbeispiel der Kontext \((u,u)\), um
\(u\not\equiv_E u^2\) zu beweisen, aber nicht allein, um
\(u^3\eqE0\) zu beweisen?

\textbf{Lösung.}
Im ersten Fall ist eine verlangte Gleichheit falsch: \(u^3\ne u^4\).
Das widerlegt die Allquantifikation. Im zweiten Fall stimmt zunächst
nur eine von 25 verlangten Gleichheiten. Erst das Argument über die
Mindestexponentensumme fünf zeigt alle übrigen Gleichheiten mit.

\textbf{3. Können in einer Halbgruppe mit Eins zwei verschiedene Elemente
in derselben Klasse liegen?}

\textbf{Lösung.}
Schon der Kontext \((1,1)\) trennt alle verschiedenen Elemente.
Alle Klassen sind Einermengen.

\textbf{4. Was heißt \enquote{ein Profil} im Rechtecksband?}
Ist bei einer einzigen Kontextklasse zwangsläufig jedes Dreierprodukt gleich?

\textbf{Lösung.}
Nein. Im Rechtecksband ergibt der Kontext die Zeile des linken und
die Spalte des rechten Faktors. Diese können wechseln. Nur die
Änderung des mittleren Faktors bleibt unbemerkt.

\clearpage
\Abschnitt{19. Woher kommt die Idee? Ein geschichtlicher Ausblick}{sec:history}

Unsere Relation gehört zu einer systematisch untersuchten Familie von
Halbgruppenkongruenzen. Samuel J. L. Kopamu führte diese Familie 1995
in seiner Arbeit \emph{On semigroup species} ein. In seinem Artikel
\emph{The Concept of Structural Regularity} von 1996 nennt er die
frühere Arbeit ausdrücklich als Ursprung der Familie; auf Seite 436
steht ihre Definition als Formel (1.1) [1, 2].

\textbf{Von einzelnen Faktoren zu ganzen Produktblöcken.}
Der Gedanke lässt sich an unserem Beispiel weiterentwickeln. Bisher
stand genau ein Faktor links und einer rechts. Man kann stattdessen
vorgeben, dass links ein Produkt aus zwei Faktoren und rechts ein
Produkt aus drei Faktoren stehen soll. Wegen der Assoziativität kommt
es nur auf die Werte dieser Produktblöcke an.

\Rand{\(n,m\); \(C^{[n]}\)}{Längen der äußeren Produktblöcke; Menge aller Werte aus genau \(n\) Faktoren.}
Für eine positive ganze Zahl \(n\) schreiben wir in diesem Ausblick
\[
  C^{[n]}=\{c_1\cdots c_n\mid c_1,\ldots,c_n\in C\}.
\]
\Rand{\(c_1,\ldots,c_n\)}{Die einzelnen Faktoren eines Produktblocks.}
Gemeint ist die Menge der \emph{Produktwerte}; verschiedene Folgen
von Faktoren können denselben Wert ergeben. Insbesondere ist
\(C^{[1]}=C\), während \(C^{[2]}\) die Menge aller Zweierprodukte ist.
Die eckigen Klammern im Exponenten unterscheiden diese Schreibweise
von \(C^1\), der Halbgruppe mit einer bei Bedarf hinzugefügten Eins
aus Abschnitt~14. Kopamu verwendet für Produktmengen die übliche
Potenzschreibweise; wir passen die Zeichen hier an unsere Lesefassung an.

Soll auf einer Seite kein Faktor stehen, benutzen wir die Eins von
\(C^1\) und setzen
\Rand{\(C^{[0]}\)}{Der fehlende äußere Block, dargestellt durch \(\{1\}\) in \(C^1\).}
\[
  C^{[0]}=\{1\}\subseteq C^1.
\]
Dies lässt nur den betreffenden Block weg. Bei einer positiven
Blocklänge stammen weiterhin alle einzelnen Faktoren aus \(C\).

Mit diesen Bezeichnungen lautet die Familie, für nichtnegative ganze
Zahlen \(n,m\) [2, 4],
\Rand{\(\theta(n,m)\)}{Kongruenz: gleiche Produktantworten bei linken Blöcken der Länge \(n\) und rechten der Länge \(m\).}
\[
  \theta(n,m)=\{(a,b)\in C\times C\mid
     xay=xby\text{ für alle }x\in C^{[n]},\ y\in C^{[m]}\}.
\]
Dabei gelten insbesondere
\[
\begin{aligned}
 (a,b)\in\theta(1,1)&\quad\Longleftrightarrow\quad a\eqE b
                         &&\text{für }E\ne\varnothing,\\
 (a,b)\in\theta(0,0)&\quad\Longleftrightarrow\quad a=b.
\end{aligned}
\]
Die erste Zeile ist genau unsere zweiseitige Relation. In der zweiten
Zeile fehlen beide äußeren Blöcke; der Test liefert unmittelbar das
mittlere Element. Die Fälle \(\theta(1,0)\) und \(\theta(0,1)\)
vergleichen dagegen nur linksseitige beziehungsweise rechtsseitige
Produktantworten.

\begin{Lesebeispiel}{Längere äußere Blöcke können mehr Unterschiede verdecken}
Im Fünferbeispiel mit \(u^5=0\) ist
\[
  C^{[2]}=\{u^2,u^3,u^4,0\}.
\]
Für \(\theta(2,2)\) wählen wir beide äußeren Blöcke aus dieser Menge.
Ist keiner von ihnen Null und auch die Mitte von Null verschieden,
beträgt die gesamte Exponentensumme mindestens \(2+1+2=5\).
Das Ergebnis ist Null. Mit einem Nullfaktor ist es ebenfalls Null.
Somit liefert jeder erlaubte Test für jedes mittlere Element denselben
Wert: \(\theta(2,2)\) hat nur eine Klasse.

Bei \(\theta(1,1)\) hatten wir noch drei Klassen. Längere äußere
Blöcke machen die Beobachtung hier also gröber. Das ist möglich, weil
Produktblöcke nicht jeden ursprünglichen Faktorwert annehmen müssen.
In diesem Beispiel kann ein Zweierblock etwa nicht den Wert \(u\) haben.
\end{Lesebeispiel}

Dieser Vergleich benutzt stets die ursprüngliche Halbgruppe und ihre
Produktwerte. Er ist vom wiederholten Quotientenbilden in Abschnitt~16
zu unterscheiden: Dort wurden bereits die Testantworten selbst zu
Klassen zusammengefasst.

\textbf{Wozu wurde die Familie untersucht?}
Kopamus Dissertation von 1996 nennt als Ziele das Untersuchen,
Konstruieren und Klassifizieren von Halbgruppen [3]. Der Zugang über
solche Kongruenzen kann Strukturen sichtbar machen, die in der
ursprünglichen Multiplikation weniger leicht zu erkennen sind.
Ein Schwerpunkt seines Artikels von 1996 ist die Frage, wann ein
entsprechender Quotient eine reguläre Halbgruppe ist [2].
Regulär bedeutet hier: Zu jedem Element \(a\) gibt es ein Element
\(b\), für das \(aba=a\) gilt. Auch eine Halbgruppe, die diese
Eigenschaft selbst nicht hat, kann einen regulären Quotienten besitzen.
So verbindet die Konstruktion allgemeinere Halbgruppen mit einer
bereits intensiv untersuchten Strukturklasse.

\textbf{Wer hat die Relation zuerst entdeckt?}
Die Quellen belegen Kopamus Einführung der \emph{gesamten Familie}
1995 und die Einordnung unserer Relation als ihren Fall
\(\theta(1,1)\). Sie belegen keine absolute Erstentdeckung dieses
einzelnen Falls. Eine einfache Konstruktion kann schon vorher in
anderen Zusammenhängen benutzt worden sein. Deshalb lautet die
historisch belegte Aussage hier:

\begin{Lesebeispiel}{Die gesicherte historische Zuordnung}
Für \(E\ne\varnothing\) ist die hier untersuchte Relation der Fall
\(\theta(1,1)\) einer von Samuel J. L. Kopamu 1995 eingeführten
Familie von Halbgruppenkongruenzen. Wer diesen einzelnen Fall erstmals
betrachtet hat, ist mit den hier ausgewerteten Quellen nicht entschieden.
\end{Lesebeispiel}

\Needspace{20\baselineskip}
\textbf{Quellen zum Weiterlesen.}
Die beiden ersten Arbeiten sind die ursprünglichen Fachartikel;
der dritte Eintrag führt zum Dissertationsnachweis mit ausführlicher
Zusammenfassung. Die geschichtliche Zuordnung stützt sich vor allem
auf die ausdrückliche Rückverweisung in [2].

\begin{enumerate}
\item Samuel J. L. Kopamu: \emph{On semigroup species}.
  Communications in Algebra \textbf{23} (1995), 5513--5537.
  \href{https://doi.org/10.1080/00927879508825550}{DOI: 10.1080/00927879508825550}.
\item Samuel J. L. Kopamu: \emph{The Concept of Structural Regularity}.
  Portugaliae Mathematica \textbf{53} (1996), 435--456.
  Besonders S.~436, Formel (1.1), und S.~455, Literaturhinweis [21].
  \href{https://emis.de/ft/52678}{Frei zugänglicher Originalartikel}.
\item Samuel Joseph Lyambian Kopamu: \emph{On a family of semigroup
  congruences}. Dissertation, University of St Andrews, 1996.
  \href{https://research-repository.st-andrews.ac.uk/handle/10023/13699}{Dissertationsnachweis und Zusammenfassung}.
\item Phillip M. Edwards, Peter M. Higgins und Samuel L. Kopamu:
  \emph{Remarks concerning Lallement's Lemma and coextensions of
  eventually regular semigroups}. Acta Scientiarum Mathematicarum
  (Szeged) \textbf{67} (2001), 105--120. Die Familie steht auf S.~106.
  \href{https://acta.hu/download.phtml?id=2598}{Frei zugänglicher Originalartikel}.
\end{enumerate}

\clearpage
\Abschnitt{20. Die wichtigsten Zeichen zum Nachschlagen}{sec:notation}

Die Randnotizen stehen bei der Einführung der Zeichen. Die folgende
Liste sammelt die dauerhaften Bezeichnungen; die Namen einzelner
Beispiele gelten jeweils nur innerhalb ihres Beispiels.

\begingroup
\renewcommand{\arraystretch}{1.3}
\begin{tabularx}{\linewidth}{@{}lX@{}}
\toprule
Zeichen&Bedeutung\\\midrule
\(C\)&Grundmenge der untersuchten Halbgruppe.\\
\(ab\), \(xay\)&Produkte in \(C\); die Klammerung ist wegen Assoziativität frei.\\
\(E\)&Parametermenge der Ausgangsformel; \(e\) kommt im Produkt nicht vor.\\
\(a\equiv_E b\)&Die Relation aus \eqref{eq:original}; für nichtleeres \(E\) Gleichheit aller Kontextantworten.\\
\((x,y)\in C\times C\)&Ein zweiseitiger Kontext mit linkem Faktor \(x\) und rechtem Faktor \(y\).\\
\(P_a\)&Kontextprofil \((x,y)\mapsto xay\).\\
\([a]_E\)&Äquivalenzklasse des Grundelements \(a\).\\
\(Q=C/{\equiv_E}\)&Menge der verschiedenen Klassen.\\
\(\ast\)&Quotientenprodukt \([a]_E\ast[b]_E=[ab]_E\).\\
\(AB\)&Mengenprodukt \(\{ab\mid a\in A,b\in B\}\).\\
\(\pi\)&Projektion \(a\mapsto[a]_E\).\\
\(\mathcal P,\rho\)&Menge vorkommender Profile und Zuordnung \(a\mapsto P_a\).\\
\(0\), \(1\)&Null- und Einselement, sofern vorhanden.\\
\(C^1\)&Halbgruppe mit einer bei Bedarf hinzugefügten Eins.\\
\(D\), \(\sim_D\)&Eingeschränkter Testvorrat und zugehörige Testrelation.\\
\(n,m\)&Nichtnegative ganze Zahlen: Längen der äußeren Produktblöcke im Ausblick.\\
\(C^{[n]}\)&Menge der Werte aus genau \(n\) Faktoren; \(C^{[0]}=\{1\}\subseteq C^1\).\\
\(\theta(n,m)\)&Kongruenz aus gleichen Antworten bei äußeren Blöcken der Längen \(n,m\).\\
\bottomrule
\end{tabularx}
\endgroup

\end{document}
